% !TEX program = pdflatex
% Depth Saturation Predicts When Looping Cannot Help
\documentclass[10pt,twocolumn,letterpaper]{article}

\usepackage[margin=1in]{geometry}
\usepackage{amsmath,amssymb}
\usepackage{graphicx}
\usepackage{booktabs}
\usepackage{xcolor}
\usepackage[colorlinks=true,linkcolor=blue,citecolor=blue,urlcolor=blue]{hyperref}
\usepackage{caption}
\usepackage{enumitem}
\usepackage{times}

\captionsetup{font=small,labelfont=bf}

\title{\vspace{-1.2em}Depth Saturation Predicts When Looping Cannot Help:\A Controlled Negative Result on Looped Transformers at 350M}
\author{Kru Zeppesey\\small Jaume Project\\small ModelScope: \texttt{LightNeko/Jaume}}
\date{}

\begin{document}
\maketitle

\begin{abstract}
Looped and recurrent transformers are claimed to buy reasoning depth without
buying parameters: apply the same block $ times and get \times$ the
effective depth. We test this claim at 350M parameters on a single Ascend
910B3 and find it does not hold here, for a structural reason rather than a
tuning failure. Our central contribution is a cheap \emph{depth-saturation
criterion}: train plain 8L and 16L models on the target task before trying a
loop. If 16L does not clearly beat 8L, the task is already computed in a
single forward pass and looping can only repeat it. We validate the criterion
across eight task structures and four modalities, including a task that is
genuinely learnable (8-digit addition, 0.89 exact match) and genuinely
serial (carry propagation). In every case 8 layers saturate. We further give
a mechanism: inter-round updates form a contraction with geometric decay
ratio $\approx 0.6$, so extra rounds converge to a fixed point rather than
reasoning further.
\end{abstract}

\section{Introduction}

A looped transformer applies the same block repeatedly, so parameters stay
fixed while effective depth grows. Recent work reports gains: LoopMTP
\cite{loopmtp}, ECHO \cite{echo}, LLA \cite{lla}. These are positive results
from large models trained for long. The natural question is when the
mechanism cannot work at all.

We ran a pre-registered equal-compute comparison at 350M and it failed:
the loop configuration matched for wall-clock was 0.99 nats worse than plain
training. Rather than tune further, we measured \emph{why}, and the
measurement turned into a predictor.

\paragraph{Contributions.}
\begin{enumerate}[leftmargin=1.2em,itemsep=1pt,topsep=2pt]
\item A cheap pre-flight criterion: depth-scan the task first
      (Section~\ref{sec:criterion}).
\item Eight task structures and four modalities, all consistent
      (Section~\ref{sec:validation}).
\item A mechanism: rounds contract geometrically, ratio $\approx 0.6$
      (Section~\ref{sec:mechanism}).
\item The distinction between \emph{not learnable} and \emph{not depth-limited},
      which a naive reading of the parity result would conflate
      (Section~\ref{sec:learnability}).
\end{enumerate}

\section{Setup}

All runs use Jaume-Causal, a 350M MoE transformer (16 layers, =512$,
4 experts, top-2) on one Ascend 910B3 with 61\,GiB HBM. Training is
AdamW, bf16 autocast, learning rate \times10^{-4}$ with cosine decay.
Every comparison uses identical data, seed, and step budget unless stated.

\textbf{Loop implementation.} The looped variant re-applies the full stack
$ times with per-round gating, following the standard full loop,
repeat position'' form. Latent-alignment supervision is available as an
ablation.

\section{The equal-compute result}
\label{sec:equal}

Pre-registered criterion: at equal wall-clock, the looped configuration C
must beat plain configuration A. It did not.

\begin{table}[t]
\centering
\small
\begin{tabular}{lcccc}
\toprule
Config & blocks & depth & val loss & wall \
\midrule
A base        & 40k & 16 & \textbf{4.9406} & 17.8m \
B loop (4x)   & 40k & 64 & 4.8975 & 44.8m \
C loop equal  & 10k & 64 & 5.9282 & 11m \
D C + align   & 10k & 64 & 5.9259 & 11m \
\bottomrule
\end{tabular}
\caption{Equal-data and equal-compute comparison. B's advantage over A is
compute, not looping: at equal wall-clock (C) the loop is 0.99 nats worse.}
\label{tab:equal}
\end{table}

Table~\ref{tab:equal} is the whole story in one table. B looks better than A
only because it consumed 2.5$\times$ the compute. Matching wall-clock (C)
reverses the ordering by nearly a nat. Alignment supervision (D) does nothing.

\section{Mechanism: rounds contract}
\label{sec:mechanism}

We probed the loop-4 checkpoint at inference with =1\ldots8$ (same weights,
no retraining). Two measurements:

\paragraph{Inter-round relative distance.}
\[
T{=}8:\; 0.552 \to 0.363 \to 0.213 \to 0.137 \to 0.097 \to 0.074 \to 0.060 .
\]
Monotone geometric decay, ratio $\approx 0.6$. Rounds 1--2 change the
representation by 55\% and 36\%; round 7$\to changes it by 6\%.

\paragraph{Principal-subspace overlap.}
\[
T{=}8:\; 0.681,\ 0.848,\ 0.945,\ 0.973,\ 0.982,\ 0.993,\ 0.991 .
\]
Overlap saturates near 0.99. All 512 dimensions are moving''
(dim-var ratio $=1.0$): the loop is not splitting labour between
accumulating stable facts and advancing reasoning. It is making
decaying updates of the same kind.

\textbf{Reading.} The loop is approximately a \emph{contraction}: each round
shrinks the update by a constant factor and converges to a fixed point.
Extra rounds are not thinking, they are converging. This explains the
equal-compute loss and also reinterprets latent overthinking'': the fixed
point is already reached.

\section{The depth-saturation criterion}
\label{sec:criterion}

If rounds are extra depth, a loop can only help when depth helps. So measure
that directly, before spending compute on the loop:

\begin{quote}
Train plain 8L and 16L models on the target task, identical data and steps.
\begin{itemize}[leftmargin=1em,itemsep=1pt,topsep=2pt]
\item 16L beats 8L by $>0.01$ (nats or task metric): the task is
      depth-limited; a loop has room.
\item 16L $\approx$ 8L: the task is computed in one pass; looping is
      structurally guaranteed to lose to more data at equal compute.
\end{itemize}
\end{quote}

The scan costs about 0.1 GPU-hour per depth on this hardware, i.e.\ under
0.3 hours total, versus 10+ hours for a full loop comparison.

\section{Validation across tasks}
\label{sec:validation}

A criterion is only useful if it is hard to break. We tried hard, across
eight task structures and four modalities.

\begin{table*}[t]
\centering
\small
\begin{tabular}{llcccl}
\toprule
Task & Intended serial structure & 8L & 16L & Learnable & Depth-limited \
\midrule
Chain lookup          & hBchop rule composition        & 0.0000 & 0.0000 & yes & no (closure in one pass) \
Register program      & 16 sequential mod-13 ops        & 0.0000 & 0.0000 & yes & no (parallel width) \
Sequence reduction    & 256-number sum/max              & 0.0000 & 0.0000 & yes & no (parallel reduce) \
Two-stage derivation  & stage 2 uses stage 1 result     & 0.1815 & 0.1812 & yes & no \
Colour counting (VLM) & count \times3$ attribute pairs & 0.983 & 0.986 & yes & no \
Constraint path       & same-chain, propagation depth $ & 0.0000 & 0.0000 & yes & no (transitive closure) \
Parity / mod-count    & running mod-2/3 state           & 0.50 & 0.56 & \textbf{no} & n/a (not learnable) \
8-digit addition      & carry chain, exact match        & \textbf{0.8900} & \textbf{0.8967} & yes & no \
\bottomrule
\end{tabular}
\caption{Depth scan across eight task structures. The last row is the
strictest: genuinely learnable, genuinely serial (carry propagation), strict
exact-match scoring. 16L gains only 0.0067 over 8L while costing 1.8$\times$
the time.}
\label{tab:tasks}
\end{table*}

\paragraph{Why the addition row matters.} Earlier rows can be explained away.
Chain lookup and constraint paths are computed by transitive closure;
sequence reduction and counting are parallel aggregation. But 8-digit
addition is the task the loop literature itself uses to argue for depth:
the carry into column $ depends on the carry out of column -1$, and the
answer digits appear nowhere in the input. It is learnable (0.89 exact
match) and strict (one wrong digit is zero). Yet 8 layers saturate: 0.8900
vs 0.8967.

The mechanism is that attention parallelises the carry. Whether column $
\emph{produces} a carry is a local function of the two input digits; a model
can read all columns at once and compute the carry propagation as a
closed-form function of the inputs, without ever passing state column by
column.

\section{Not learnable is not the same as not depth-limited}
\label{sec:learnability}

Parity is the standard hard task for shallow models, so it looks like the
ideal stress test. It is not, and the distinction matters.

At 8, 16 and 32 layers parity stays at chance (0.50--0.58) with
cross-entropy not decreasing. This is \emph{optimisation invisibility}, not a
depth limit: the gradient signal is too weak for the function to be learned
at any depth. Reading it as not enough depth'' would produce exactly the
wrong conclusion.

A usable criterion therefore has two clauses, not one:
\textbf{(a) the task must be learnable, (b) 16L must beat 8L.} Only tasks
satisfying both are candidates for looping. The addition experiment is the
one that satisfies both, and it still saturates.

\section{Implications}

\begin{enumerate}[leftmargin=1.2em,itemsep=2pt,topsep=2pt]
\item \textbf{Screen before you loop.} A 0.3-hour depth scan can rule out a
      direction that would otherwise cost 10+ hours.
\item \textbf{Scaling this model up will not rescue looping.} If 8 layers
      already saturate, more layers saturate sooner, and loops have less to
      add, not more.
\item \textbf{Look for tasks where attention cannot collapse the serial
      structure} --- incompressible state that must be carried, not
      recomputed: game-tree search with bounds, or an interactive
      environment with memory. These remain open.
\end{enumerate}

\section{Limitations}

One model size (350M), one architecture family, largely synthetic tasks.
The criterion is validated as a \emph{negative} screen: when 8L saturates, the
loop is ruled out. We do not claim the converse --- a depth-limited task is
necessary for looping to help, not sufficient. The positive direction (what
happens when 16L does beat 8L) is untested here because we could not
construct such a task at this scale.

\section{Conclusion}

At 350M, on every task structure and modality we could construct, eight
layers suffice. Where that holds, looping cannot pay, and a cheap depth scan
says so in advance. The transferable result is the scan itself: measure
whether depth helps before paying for a loop.

\bibliographystyle{plain}
\bibliography{refs}

\end{document}
